% sections/07_conclusion.tex
% 第7章：综合判断

\section{综合判断}

本章将前六章的分析结果按认知层级收束为三个明确类别：已知事实（无需进一步验证即可确认）、可计算链（在当前参数集和假设下可给出量化范围）、当前不可判定（需要更多数据才能给出可靠判断）。本章是全文的逻辑终点，不引入新证据或新计算。

\subsection{已知事实}

以下结论基于物理定律或可公开验证的权威资料，不需要进一步验证即可确认：

\begin{enumerate}[leftmargin=*]
    \item \textbf{斯特藩-玻尔兹曼定律不否定 1\;MW 级轨道散热。} AI1 的 1,400\;W/m\textsuperscript{2}（双面辐射、$\varepsilon\approx0.90$、$T\approx342.5$\;K）在物理上自洽且可由当前材料实现。但该工作点几乎用尽了高辐射率涂层 EOL 退化后的可用裕度——物理可行不等于裕度充足。（来自：第4章 S-B 推导）

    \item \textbf{散热密度的翻倍（$\to$2,800+\;W/m\textsuperscript{2}）在 2035 年前无单一技术路径可支撑。} 组合路径（高温芯片 + 耐久涂层 + 可展开结构）可实现 40--50\% 改善。散热密度是物理定律决定的硬约束，不是可以通过"工程规模"无限压降的变量。（来自：第4章散热密度边界分析）

    \item \textbf{LEO 真实热环境不是"理想冷库"。} 太阳直射、地球反照和地球红外辐射使有效散热能力比理想真空低 15--40\%。轨道昼夜食影循环（每 96.7 分钟一次，食影 35.5 分钟）引入了不可消除的储能约束。（来自：第4章 LEO 热环境分析）

    \item \textbf{发射成本在过去 15 年下降了约 100 倍}，从 Shuttle 约 \$54,500/kg 到 Falcon 9 复用约 \$2,720/kg。Starship V3 的目标是 \$100--200/kg（全复用模式），但截至 2026 年 6 月尚未实现全复用\cite{nasalaunch2018,faa2025boca,faa2026ksc}。（来自：第2章发射单价口径说明）

    \item \textbf{在当前公开的 AI1 翼展约束和现有展开技术先例下，HJT 路线显著不利。} AI1 几乎必然使用 GaAs（或等价 III-V 族）太阳电池。这一推论的置信度极高——625\;m\textsuperscript{2} 独立物理推导与 70\;m 翼展在设计参数上精确自洽（误差 <1\%），构成强交叉验证。（来自：第5章翼展约束分析）
\end{enumerate}

\subsection{可计算链}

以下结论在当前参数集和假设下可给出量化范围。其精度受限于部分参数的四档推断性质（14/47 参数），但方向性判断稳健：

\begin{enumerate}[leftmargin=*]
    \item \textbf{10年-GaAs 主链 TCO 约\$0.77B（区间\$0.55B--\$1.1B，取决于关键假设变动）}（1\;MW 持续 IT / 10 年总窗口 / \$200/kg）。这是当前证据最充分的单代主链基线结果。光伏阵列（约 50.4\%）和系统裕量/冗余（约 20.0\%）是总成本的最大两项，合计占约 70.4\%。发射成本仅占约 1.8\%。（来自：第6章 TCO 拆解）

    \item \textbf{5年-GaAs 对照 TCO 约 \$1,341.26M}，比主链贵约 75\%。寿命缩短至 5 年导致需重射一次，制造成本近似翻倍。\textbf{在轨寿命是比发射成本重要得多的经济杠杆}——将寿命从 5 年延长至 10 年省下的成本（\$576M）相当于将发射单价从 \$200/kg 降至约 \$0/kg 的效果。（来自：第6章对照分支对比）

    \item \textbf{5年-HJT 补齐估计 TCO 约 \$1,005M}，在成本上不优于 10年-GaAs。在当前公开的 AI1 翼展约束和现有展开技术先例下，HJT 路线显著不利。（来自：第6章 HJT 并行路线分析）

    \item \textbf{单星质量基准约 8.48\;t}（10年-GaAs），含 9 个可独立追溯的子系统。计算载荷质量约 1,500\;kg，证据链从此前的"B1 二手转述"提升至"A1 地面硬件锚点（GB300 NVL72）+ 三项可验证的空间化增量"。（来自：第5章 bottom-up 质量链）

    \item \textbf{发射单价 \$200/kg 敏感性}：从 \$200/kg 降至 \$100/kg，总成本仅下降约 0.9\%。降至 \$50/kg，总成本仍约 \$751M。发射成本对总 TCO 的影响有限——这是本报告最稳健的定量结论之一。（来自：第6章发射单价敏感性分析）

    \item \textbf{马斯克的 300--500\;GW/年部署声明与当前 FAA 批准的年发射上限（69 次）存在约 3 个数量级的偏差。}即使在 FAA 上限大幅提升的乐观情景下，闭合这一差距需要运力、频次和单星质量三个参数的组合突变。（来自：第5章发射运力现实检验）
\end{enumerate}

\subsection{当前不可判定}

以下问题在 2026 年中期的公开证据水平下无法给出可靠判断：

\begin{enumerate}[leftmargin=*]
    \item \textbf{轨道 GPU 实际故障率}。当前估计为 5--15\%/年（四档推断），需要 Rubin Space-1 在轨数据才能闭合为更精确的值。故障率的真实值将直接影响重射间隔、冗余设计和 TCO。（与残余风险 R3 关联）

    \item \textbf{多参数同步改善的综合概率}。太空算力的商业竞争力需要光伏成本下降、在轨寿命延长、AIT 成本压缩和发射成本下降等在\emph{同一时间窗口}内同步发生。每个参数的单独改善概率可以估计，但四参数同步改善的概率无法精确量化。本报告判断该概率为低——基于"同步性"本身的概率折扣（$P(A \cap B \cap C \cap D) \ll \min(P(A), P(B), P(C), P(D))$）和四个参数中每一项距离可竞争水平的当前差距。

    \item \textbf{HJT 空间级阵列系统级采购价格的完整闭合}。价格链有三个独立锚点（地面 HJT 批发价、东方日升空间级报价、Starlink 级阵列系统参考），可以给出区间估计（第6章已给出补齐估计约 \$15M/颗阵列系统，总成本约 \$1,005M），但空间级系统采购的一手公开报价仍然是缺口。然而，即使此缺口闭合，也不改变"5年-HJT 在 AI1 约束下不优于 10年-GaAs"的结论。（与残余风险 R4 关联）

    \item \textbf{AI1 一手总质量和收拢尺寸}。虽已通过 bottom-up 质量链部分缓解（地面硬件锚点 + 空间化增量），但无法闭合到 SpaceX 的一手官方总质量。8.48\;t 是当前最佳估计，不是已验证事实。（与残余风险 R1 关联）

    \item \textbf{Rubin 化后的整星 BOM 和系统级改善程度}。NVIDIA 已发布 Rubin 的 GPU 级提升（推理性能比 H100 提升 25 倍），但整星级别的 BOM 变化（包括辐射加固、热管理、供电等系统级适配）缺乏一手公开数据。Rubin 会缩小太空与地面之间的差距，但缩小到何种程度当前无法判断。（与残余风险 R3 关联）
\end{enumerate}

以上不可判定项中，第1、2项是方向性判断的核心不确定性来源；第3、4、5项是定量精度的不确定性来源，但不改变方向性结论。

\subsection{时间线展望}

基于"物理可行 $\to$ 工程脆弱 $\to$ 商业远距"的三层递进判断，本报告对太空算力的发展时间线给出以下展望：

\begin{figure}[H]
\centering
% 结论章时间线图 — 第7章用
% 三阶段前瞻判断时间线：原型验证 → 有限部署 → 条件性扩张
% 明确为"前瞻判断时间线"，非历史事件回顾
\begin{tikzpicture}[
  phase/.style={rectangle, draw=gray!70, rounded corners, minimum width=3.0cm,
                minimum height=1.2cm, align=center, font=\small},
  judge/.style={text width=2.8cm, align=center, font=\footnotesize, text=black!70},
  arrow/.style={thick, draw=gray!60, ->, >=stealth},
  note/.style={font=\small\itshape, text=gray!50!black, align=center}
]
  % ===== 底部时间轴 =====
  \draw[thick, ->, >=stealth] (-0.3,0) -- (10,0);

  % 时间刻度
  \foreach \x/\l in {0/2026, 2.8/2028, 5.6/2030, 8.4/2035}
    \draw (\x,0.08) -- (\x,-0.08) node[below, font=\footnotesize] {\l};

  % ===== 阶段一：2026--2028 原型验证期 =====
  \node[phase, fill=gray!20] (p1) at (1.4,0.9) {\textbf{2026--2028}\\[2pt]\textbf{原型验证期}};
  \node[judge] at (1.4,-0.9) {经济性不可行\\技术验证有意义};

  % ===== 阶段二：2028--2030 有限部署期 =====
  \node[phase, fill=cyan!18!gray!14] (p2) at (4.2,0.9) {\textbf{2028--2030}\\[2pt]\textbf{有限部署期}};
  \node[judge] at (4.2,-0.9) {利基场景可行\\规模化不可行};

  % ===== 阶段三：2030--2035 条件性扩张期 =====
  \node[phase, fill=blue!16!gray!20] (p3) at (7.0,0.9) {\textbf{2030--2035}\\[2pt]\textbf{条件性扩张期}};
  \node[judge] at (7.0,-0.9) {取决于多参数\\同步改善概率};

  % ===== 阶段间演进箭头 =====
  \draw[arrow] (p1.east) -- (p2.west);
  \draw[arrow] (p2.east) -- (p3.west);

  % ===== 底部说明 =====
  \node[note] at (4.6,-2.0) {本图为前瞻判断时间线，基于第4--6章分析结论，不代表实际发生预测};
\end{tikzpicture}

\caption{太空算力三阶段前瞻判断时间线——从原型验证到条件性扩张的演进路径。每阶段的核心判断基于第4--6章的物理、工程与商业分析结论。}
\label{fig:timeline_outlook}
\end{figure}

\begin{enumerate}[leftmargin=*]
    \item \textbf{2026--2028：原型验证期}。AI1 首批 2 颗原型（计划 2027 年发射）和 Rubin Space-1 在轨验证数据是该阶段最重要的信息增量。此阶段的商业可行性为零——原型阶段的成本不能与地面量产数据中心的成本进行比较。关键在于验证：轨道 GPU 实际故障率、散热系统在真实 LEO 环境下的性能、光伏阵列在轨退化率。\emph{判定：经济性不可行，技术验证有意义。}

    \item \textbf{2028--2030：有限部署期}。如果原型验证结果积极，太空算力可能在 2028--2030 年间进入有限部署阶段——军事/情报场景（对物理安全、全球覆盖和不可拦截性有溢价支付意愿）和低延迟推理利基场景（利用 LEO 全球覆盖 <20--30 ms 延迟优势）可能是首批商业客户。此阶段的部署规模预计为数十至数百颗，而非百万颗。\emph{判定：利基场景可行，规模化不可行。}

    \item \textbf{2030--2035：条件性扩张期}。如果以下四个条件同步达标——(1) 在轨寿命稳定 $\geq$ 7 年；(2) 光伏阵列成本降至 <\$50/W\_BOL；(3) AIT 成本压缩至 Starlink 量产级；(4) 发射成本 <\$100/kg——太空算力可能开始向规模化商业应用扩张。四个条件的同步概率被评估为 <25\%，因此 2030--2035 条件性扩张是"可能但不确定"而非"大概率发生"。\emph{判定：取决于多参数同步改善概率，\\
    当前不足以支撑明确承诺。}
\end{enumerate}

\subsection{结论的稳健性说明}

\begin{itemize}[leftmargin=*]
    \item \textbf{物理结论（第4章）高度稳健}：仅依赖物理常数（$\sigma = 5.6704\times10^{-8}$\;W\;m\textsuperscript{-2}\;K\textsuperscript{-4}）和可公开验证的物理定律，不依赖任何公司披露或商业假设。任何未来的技术进展或信息增量都不会改变斯特藩-玻尔兹曼定律和 LEO 轨道力学。
    
    \item \textbf{工程结论（第5章）中等稳健}：bottom-up 质量链的"地基"——GB300 NVL72 等地面硬件规格——来自二级权威资料（Lenovo/NVIDIA 官方），锚点可靠。三项空间化增量中的辐射防护和热管理改造均有独立文献锚点（三档），但结构改造增量仍为行业经验估计（四档）。整星总质量的 9 子系统拆解中约 35\% 的成本项涉及四档推断——如果 AI1 的一手总质量或收拢尺寸公开，整星质量的置信区间将显著收窄。
    
    \item \textbf{商业结论（第6章）对参数假设敏感}：TCO 模型中 14/47 参数为四档合理推断。核心不确定性不在于任何单一参数，而在于多参数同步改善的概率。换言之，决定商业可行性的不是某一项技术单独突破，而是多个关键参数能否在同一时间窗口内一起改善。

    \item \textbf{研发非经常性工程费用（NRE）未计入 TCO 模型。} 百万颗量产后可摊薄至可忽略水平，但在早期 1--100\;MW 验证阶段不可忽略。NRE 主要来源包括：卫星平台 GNC/热控/电源/结构的非重复设计验证、地面段定制化开发、批量产线一次性投资。

    \item \textbf{方法论稳健性}：本报告从物理定律（斯特藩-玻尔兹曼）出发，经工程硬件锚点（NVIDIA GB300 NVL72），到独立 Python 建模的完整计算链，遵循第一性原理自底向上验证的方法论。主结论尽量追溯到公开可独立核验的原始来源；光伏阵列（$\sim$50\%）已通过二档/三档交叉锚定，确保方向性结论不会因四档参数的合理变动而发生数量级翻转。
\end{itemize}

\subsection{分析立场与方法}

本报告以场景、口径和证据三层分层处理争议问题：

\begin{itemize}[leftmargin=*]
    \item \textbf{场景分层}：TCO 必须与具体假设一同阅读。本报告给出了三个可比分支的完整结果，并标注了各自边界条件。
    \item \textbf{口径分层}：不同口径下的数字不能直接比较。附录 A 提供了 7 个场景 $\times$ 9 个维度的完整对齐矩阵。
    \item \textbf{证据分层}：47 个关键参数全部按四档证据体系分级并裁决。三档和四档参数用于表述"当前证据水平下的最佳估计"，而不直接写成已证实事实。
    \item \textbf{保留不可判定项}：对当前公开信息无法闭合的问题，报告保持"当前不可判定"这一状态，并将其与可计算链、已知事实分开处理。
\end{itemize}

如果本报告的核心判断可以用一句话概括：\textbf{太空算力的物理基础是坚实的，工程路径是可辨识的，商业前景是远距的。} 当前公开信息支持继续观察而非提前给出强商业承诺。AI1 和 Rubin Space-1 的飞行验证数据，将是下一个最重要的信息增量。

（全部可计算链数字的完整公式与代入值见附录 C；全部争议参数的去向与裁决记录见附录 B；已废弃推导链与禁线清单见附录 E；从本报告结论到原始来源的可执行追溯路径见附录 D。）

\endinput
